\section{Limitations}
\label{sec:limitations}

A paper arguing that a published conclusion is an artefact of undisclosed choices is
held to its own standard, so the accounting below is complete rather than selective.
(The original's runs to two paragraphs in its conclusion and acknowledges two items: the
cross-sectional design, and that ESGSI measures could be more sophisticated, p.~10.)
Table~\ref{tab:limitations-origin} attributes each limitation to our design or to the
index we re-implement, descriptively only. Inheritance is not a defence: an inherited
defect constrains our results exactly as much as one of our own, and the first two
items apply to \citet{lagasio2024}, and to any future user of this index family, with
equal force.

\subsection{The threshold is definitional, and the index has no cardinal scale}

Under score normalisation by $z$-score, $\ESGSI = \Z(\SEN) - \Z(\SUS)$ is a difference
of two mean-zero variables, so the rule $\ESGSI > 0 \Rightarrow$ ``Potential
ESG-washing'' partitions any sample at approximately its median by construction. Our
171 flagged documents of 343 are arithmetic, not an estimate: no figure in this paper
estimates how much ESG-washing this corpus contains, and no reader should take one as
such. The threshold is inherited. The original adopts the same zero cut with no
calibration against any external case and never reports the share it flags, although
its own published quartiles bound that share at no more than 25\,\% of 749
(\S\ref{sec:original}).

Both components are rescaled across the sample, so neither they nor the index carry
cardinal meaning. An $\ESGSI$ of $+1.5$ is a position in this corpus's distribution,
not a quantity of anything; values are not comparable across corpora, across
vocabularies, or between this study and the original. Only within-sample ordering, and
change over time with composition held fixed, are interpretable.

\subsection{Corpus dependence}

$\sustfidf$ and $\suslag$ estimate document frequencies on the analysed sample. Adding
documents changes every idf and hence every historical score, so a report's 2018 value
is not stable under corpus extension --- a live defect for a longitudinal design,
which must re-score the whole panel whenever it grows. $\susdens$ estimates no document
frequencies, so its substance component is stable under extension; that does not lift
the limitation from the index. Score normalisation is cross-sectional under every
specification, so every $\ESGSI$ reported here, including the whole yearly series of
\S\ref{sec:temporal}, would move if the panel were extended. The series is one of
positions within this corpus, not of levels.

\subsection{Confounds and unjustified constants}

Country and document type are near-perfectly confounded (\S\ref{sec:data}): 99 of 101
URDs are French, all 98 German documents are annual reports, and all seven standalone
sustainability reports are French. No cross-country or cross-type comparison in this
corpus is interpretable and we make none; the country and type means in
\S\ref{sec:data} identify neither effect. The within-Germany analysis (\S\ref{sec:temporal})
holds both fixed --- one country, one document type, 98 documents, slope $-0.301$/yr,
$p = 1.2 \times 10^{-6}$ --- but it does so on 14 firms, and it establishes that the
decline is not produced by composition, not that the pooled magnitude is right. The
cross-section is 49 firms throughout, observed seven times each, so the $p$-values in
\S\ref{sec:temporal} are anti-conservative for the reasons given there; we report no
clustered or panel-corrected inference.

The extended-index weights $w_q = w_h = 0.5$ were chosen without justification and never
varied. We report no result that depends on $\ESGSIext$, which is why no sensitivity
analysis was run --- not because the choice is defensible, and any use of $\ESGSIext$
beyond this paper would need one.

\subsection{\QUANT{} is not independent of what it adjudicates}

\QUANT{} is the criterion on which the specification choice rests. The selection of
$\susdens$ over $\suslag$ in \S\ref{sec:validity} turns on their agreement with \QUANT{}, and
the argument is a convergent-validity argument: it is worth what the independence of the two
instruments is worth. That independence fails in two ways, one of which we can remove by
recomputation and one of which we cannot. Fourteen entries of the 154-term ESG vocabulary
are also counted by \QUANT{} --- eleven by its \texttt{frameworks} group and three
(\emph{co2}, \emph{co2e}, \emph{ghg}) by its \texttt{units} group --- and they carry
7.50\,\% of all \SUS{} occurrences. The two groups that contain the shared terms are
together 36.77\,\% of all \QUANT{} hits (\texttt{frameworks} 17.22\,\%, \texttt{units}
19.55\,\%), so the contamination is not confined to a corner of the criterion. Nor is
it stationary: the \texttt{frameworks} share rises from 10.15\,\% of \QUANT{} hits in
2018 to 21.87\,\% in 2024, so the shared component grows across exactly the window over
which \S\ref{sec:temporal} reports a trend.

Two consequences follow, and they are different in kind. The first is settled by
recomputation rather than by argument, and \S\ref{sec:validity-disjoint} reports it: stripping
the 14 shared entries from the vocabulary and \QUANT{} of the \texttt{frameworks} and
\texttt{units} groups drops every correlation by roughly two fifths --- $\susdens$ from
0.6663 to 0.4112, $\sustfidf$ from 0.6817 to 0.4125, $\suslag$ from 0.3217 to 0.1799. So the
levels reported there are upper bounds and the true level of convergent validity is moderate,
not strong; that limitation stands and we state it in those terms. The ranking is not
similarly affected --- the length-normalised specifications stay roughly twice as close to
\QUANT{} as $\suslag$, and the ratio widens slightly, 2.29$\times$ against 2.07$\times$. We
had previously nothing better than an argument for why the ranking should survive, and that
argument does not hold (\S\ref{sec:validity}); what supports it is the recomputation and
nothing else. The second consequence touches the extended index: $\ESGSI$ itself does not
contain \QUANT{} and is unaffected, but $\ESGSIext$ subtracts both $\Z(\QUANT)$ and
$\Z(\SUS)$, so there the shared mass reinforces rather than cancels and the growing
\texttt{frameworks} share inflates the extended index's temporal slope.

What recomputation does not touch is the more serious of the two failures, and it is the one
that remains the most serious limitation in this paper. \QUANT{} and \SUS{} are computed over
the same extracted zones, with denominators that are two measures of the same document's
length (\S\ref{sec:validity}). Deleting shared strings does not make the two instruments
independent samples of anything; both still see only the text \S\ref{sec:extraction} admitted,
and no vocabulary edit can change that. The convergent-validity claim is therefore conditional
on the extraction stage under either set of bases.

\subsection{\QUANT{} is also crude, and shares no denominator with \SUS{}}

Independence of construction was bought at the price of commensurability. \SUS{},
\HEDGE{} and $\BREADTH$ are densities over 14.4M lemmatised, stopword-filtered tokens;
\QUANT{} is a density over 26.2M raw tokens, a base roughly 1.8 times larger and
differently composed. No two components of the index share a denominator. Score
normalisation makes them addable, so the arithmetic of $\ESGSIext$ is well defined, but
it is the score normalisation that makes them comparable rather than anything about the
quantities themselves.

The criterion is also noisy in a way we can name but have not corrected. Its
\texttt{large\_numbers} family matches any run of four or more digits with calendar
years excluded, which counts page numbers, note references, ISINs, share counts and
monetary figures alongside quantified ESG disclosure; it is 6.92\,\% of all \QUANT{}
hits (13{,}518 matches) and no part of it was validated against ESG content. And the
measure has a hard floor that one document reaches: a single report of the 343 records
zero \QUANT{} matches across all four families, so its $\Z(\QUANT)$ is bounded below by
the sample and is not a measurement of anything. That document is also among the
shortest extractions in the corpus, which is the dependence \S\ref{sec:validity} already
notes --- \QUANT{} and \SUS{} are computed over the same zones, and a report from which
extraction recovered little text is low on both for reasons unrelated to disclosure.

\subsection{The vocabulary: a recall floor and its own construction artefacts}

Recall is bounded by the vocabulary at both the extraction and counting stages: an
unlisted term is not under-weighted but absent. An external audit of 880 zones from
four FY2023 reports found 147 ESG terms present in the text and absent from our
vocabulary, over 2,970 occurrences, led by \emph{compliance} (436) and \emph{social}
(275). None of the 147 is incorporated into the adopted vocabulary, every figure in this
paper is computed on the 154 entries, and recall is bounded by that omission. The audit
shares the bound it measures: it matched lexically, inside already-extracted zones, so
paraphrase is invisible to it too.

A sector-bias analysis has since been run over all 343 reports, and it converts one item
on this list from an unresolved omission into a design choice we disclose rather than
defend as necessary. Of the audited candidates, 22 are concentrated enough in a single
industry to be called sectoral, and they are excluded; they carry 2.0\,\% of the 143
candidates' occurrence mass and 0.37\,\% of the substance numerator they would be added
to. The 22 natural-form entries yield 19 live TF-IDF entries: \emph{b corp},
\emph{iso 37001} and \emph{cop 28} do not survive preprocessing, and two of the survivors
are lemmatisation artefacts of the class this subsection goes on to describe ---
\emph{better cotton} becomes \emph{well cotton} and \emph{living income} becomes
\emph{live income}, the latter scoring two occurrences in one firm. Added on top of the
expanded 289-term robustness vocabulary of \S\ref{sec:robustness-lexicon}, they move mean
$\susdens$ from 6.5162 to 6.5406, the
flagged count from 164/343 to 167/343 and the trend slope from $-0.194$ to $-0.193$/yr;
the two indices correlate 0.99980 and three labels of 343 change.

Some of them do what the objection predicts, and that is why they are out: mean
$\susdens$ rises $0.354$ for Danone (\emph{nutrition}, \emph{food safety}), $0.111$ for
Enel (\emph{energy storage}) and $0.099$ for Schneider Electric, against a corpus mean of
$0.024$, and 34 of the 49 firms gain less than $0.02$ --- firms credited for the industry
they operate in rather than for what they disclosed. The exclusion avoids a bias we can
name and measure, and we prefer it on that ground; what it does not do is carry any
result.

The bin is not one kind of term. The rule sorts on concentration without distinguishing
its causes, so alongside industry vocabulary it holds house-style governance language ---
\emph{executive compensation} appears in 29 of the 49 companies and is excluded on the
second limb of the rule, that half its occurrences sit in one firm --- and one conference
name. Two entries were measured as zero by a tokeniser that split on letters only:
re-measured with word boundaries, \emph{sf6} is correctly sectoral and \emph{cop28} meets
the acceptance rule and sits in the sectoral file in error. And the dispersion of
\emph{board oversight} depends on the base: 8 companies as a natural-form phrase, 30 as
the lemmatised bigram the vocabulary actually matches, because stopword removal
manufactures the adjacency. The rule was applied consistently; its inputs were not all
measured on the representation the instrument scores.

Against the 96 labels of 343 that vector normalisation moves (\S\ref{sec:robustness}),
the 22 sectoral terms move three, so no claim in this paper would read differently had we
included them. The decision is disclosed and immaterial, not necessary; the note in
\texttt{esg\_terms.txt} deferring it to a sector-bias analysis overstated what such an
analysis could settle.

Precision is bounded too, by the construction rule \S\ref{sec:lexicon} defends. Passing
every vocabulary entry through the corpus preprocessor guarantees that vocabulary and
text share a representation, and it produces entries that are lemmatisation artefacts
rather than terms: \emph{esr} from \emph{esrs}, \emph{whistleblowe} from
\emph{whistleblowing}, \emph{ecovadi}, \emph{sustainalytic}, \emph{live wage},
\emph{employee wellbee}, \emph{datum breach}. The rule makes these harmless for matching,
since the corpus passes through the same transformation, but not harmless for
weighting. \emph{esr}
occurs 6{,}282 times across 74 of the 343 documents at $\mathrm{idf} = 2.52$, placing it
in the high-idf group that \S\ref{sec:measurement} identifies as the one the idf
specifications most heavily up-weight; it is the largest single entry in that group.
The rare-term critique we direct at $\suslag$ therefore applies to our own vocabulary,
and it applies to $\sustfidf$ as well as to $\suslag$. Only $\susdens$, which weights
every entry equally, is insensitive to it. We did not audit the vocabulary for
artefacts before running the analysis, and no entry was removed on that ground.

\subsection{One extraction pipeline, one tone dictionary}

Every result is conditional on the zone extraction of \S\ref{sec:extraction}, whose free
parameters --- a density threshold of 1.0 keyword per 100 words, a $\pm 1$-paragraph
context window, a 150-character minimum zone --- were set by judgement and never
varied. That stage discards roughly two thirds of each report and sits upstream of
every number here, and it has no counterpart in the original, so the limitation is ours
alone.

The extraction stage shares its vocabulary with the measurement stage, so building the
robustness vocabulary of \S\ref{sec:robustness-lexicon} varied it incidentally, and the
effect can be read off. Holding the adopted 154-term scoring vocabulary fixed and
swapping only the extraction --- 154 natural-form terms against 296 --- takes the corpus
from 14.4 to 16.4 million lemmatised tokens, lowers mean $\susdens$ from 5.897 to 5.168,
moves 7 of 343 labels and leaves the two indices correlated at 0.988. The added zones are
therefore poorer in adopted-vocabulary terms than the zones already admitted, which is
what the density threshold predicts of text entering at the extensive margin. This is one
point, not a sensitivity analysis: the three free parameters of \S\ref{sec:extraction} ---
the density threshold, the context window and the minimum zone length --- are still never
varied.

\SEN{} is a dictionary count and inherits the known limitations of
\citet{loughran2011} on non-financial narrative text: no negation or intensifier
handling, a 140-positive against 893-negative asymmetry in the stem lists, and a
bounded polarity ratio that is insensitive to the volume of matched terms --- a report
with 4 positive and 2 negative matches scores exactly as one with 400 and 200. Our
input is lemmatised text that the LM tokeniser then stems again. Substituting this
dictionary for the original's is itself a divergence (\S\ref{sec:scope}) and one more
reason the two indices are not the same instrument --- as is our vocabulary, which
cannot be aligned with the original's because the original's is unpublished
(\S\ref{sec:lexicon}).

\subsection{No external validation, and no blind specification choice}

The check in \S\ref{sec:validity} is a construct check against an internal yardstick, not
a criterion check. Neither this study nor the original tests $\ESGSI$ against an
external greenwashing benchmark --- regulatory findings, controversy data, or
hand-labelled cases. We identify that absence in the original (\S\ref{sec:original}) and
concede it here in the same terms: $\ESGSI$ is a measure whose validity rests on its
definition. The concession runs further than the original's does. Our one positive
methodological claim --- that $\susdens$ is the specification to prefer --- is decided
by a criterion we built ourselves, from our own regexes over our own extracted zones.
The original's specification is undisclosed and ours is disclosed, but neither is
anchored to anything outside its author's pipeline.

Nor was the comparison blind. The vocabulary, the extraction parameters, the \QUANT{}
regex families and the specification set were all fixed by us with the results
visible, and no part of the corpus was held out. Every individual choice is disclosed
in \S\ref{sec:repro} and each is defended where it is made, which is the disclosure
standard we hold the original to; disclosure is not the same as pre-commitment, and we
claim only the first. The classification thresholds of \S\ref{sec:robustness-lexicon} ---
twelve companies, half the occurrences in one firm --- were set later than the rest and
with the results they adjudicate already computable, so they fall inside this concession
rather than outside it.

\begin{table}[htbp]
\centering
\caption{Attribution of each limitation. ``Inherited'' means the property follows from
the index architecture of \citet{lagasio2024} and holds for the original as well.}
\label{tab:limitations-origin}
\small
\begin{tabular}{@{}p{8.4cm}l@{}}
\toprule
Limitation & Origin \\
\midrule
Zero threshold is definitional, not calibrated & Inherited \\
No cardinal scale; scores are sample-relative & Inherited \\
idf corpus dependence of $\sustfidf$, $\suslag$ & Inherited \\
Vocabulary not comparable to the original's & Inherited (list unpublished) \\
No external greenwashing benchmark & Shared \\
Country / document-type confound & Ours (corpus design) \\
Cross-section of 49 firms; no clustered inference & Ours \\
$\ESGSIext$ weights $w_q = w_h = 0.5$ & Ours \\
Criterion \QUANT{} shares 14 terms with \SUS{} (7.50\,\% of \SUS{} mass); levels in
\S\ref{sec:validity} are upper bounds & Ours \\
Criterion and measure read the same extracted zones; not removable & Ours \\
No component shares a token base with any other & Ours \\
\texttt{large\_numbers} counts non-ESG figures (6.92\,\% of \QUANT{}) & Ours \\
\QUANT{} floors at zero for 1 of 343 documents & Ours \\
Lexical recall floor (147 audited omissions) & Ours \\
Sectoral vocabulary excluded by choice (moves 3 labels of 343) & Ours (disclosed) \\
Lemmatisation artefacts up-weighted by idf (\emph{esr}) & Ours \\
Single untested extraction stage & Ours (no counterpart) \\
Extraction vocabulary varied once, incidentally (7 labels of 343) & Ours \\
\SEN{} dictionary limitations & Ours (dictionary substituted) \\
Specification chosen non-blind, no held-out data & Ours \\
\bottomrule
\end{tabular}
\end{table}
